% ═══════════════════════════════════════════════════════════
% 第 2 章：相关工作
% ═══════════════════════════════════════════════════════════
\section{相关工作}

本文的研究横跨 AI 音乐生成、可控内容生成和音乐结构化描述三个领域。
我们在本节分别综述这三个方向的核心进展，并在每小节的结尾阐明本文与已有工作的本质差异。

% ─────────────────────────────────────────────────────────
\subsection{AI 音乐生成模型}

自 AudioCraft\cite{musicgen} 的开源发布以来，基于 Transformer 的自回归音乐生成模型已成为该领域的主流范式。
MusicGen 将 EnCodec 音频分词器与 T5 文本编码器结合，通过交叉注意力机制实现文本到音频的端到端生成，
在主观听感测试中显著优于当时的扩散模型方案。
后续工作沿着两个方向推进：一是模型规模的扩展（如 ACE-Step v1.5\cite{acestep} 达到商业级生成质量），
二是生成架构的多样化——Stable Audio\cite{stableaudio} 采用 DiT 扩散模型实现了高保真的固定时长生成。

然而，上述模型的共同设计哲学是以\textbf{文本描述}为唯一控制接口。模型内部的交叉注意力层仅处理自然语言嵌入，
缺乏对结构化音乐参数的显式建模。这使得用户在尝试精确控制和弦进行、动态曲线或奏法切换时，
只能依赖模型对自然语言的隐式理解，控制精度受限于 T5 编码器的语义泛化边界。

\textbf{与本文的差异：}Camera-Music ControlNet 在 MusicGen 的交叉注意力层之上扩展了第二条控制通路——
将六维结构化 Token 编码与文本编码在序列维度直接拼接。该方案保持了对自然语言的完全兼容性，
同时以仅 0.67\% 的额外参数为代价，引入了显式的多维音乐参数控制。

% ─────────────────────────────────────────────────────────
\subsection{可控音乐生成技术}

可控内容生成的研究可追溯至 ControlNet\cite{controlnet}。
该工作通过在冻结的扩散模型旁路中添加可训练的零卷积层，实现了对边缘图、深度图、
人体姿态等多模态条件的精准响应。其核心范式——“冻结主干、旁路注入”——
已成为可控生成领域的事实标准。

在音频领域，Music ControlNet\cite{musiccontrolnet} 首次将这一范式引入音乐生成，
通过从训练音频中提取旋律、动态和节奏三种时变控制信号（以频谱图形式），
fine-tune 了基于扩散模型的音乐生成器。
相比于 MusicGen 的全文本方案，Music ControlNet 在旋律保真度上提升了 49\%，
同时将训练参数和数据需求分别降低了 35 倍和 11 倍。
然而，其控制条件停留在\textbf{频谱图特征级别}——melody spectrogram 和 envelope 曲线——
是一种非人类可读的中间表示，用户无法直接编写或修改。

MusiConGen\cite{musecong} 和 MuseControlLite\cite{muselite} 将条件类型扩展至
和弦序列和 BPM 等符号化参数，MuseControlLite 进一步支持了 melody、dynamics、
rhythm 和 audio inpainting 四种条件的任意组合，在旋律准确度上比 ControlNet 提高了 7.6\%。
但其控制信号仍然是 feature-level 的音频特征，未能提供面向用户的描述语言。

\textbf{与本文的差异：}相较于 Music ControlNet 的频谱图级条件和 MusiConGen 的单维度和弦控制，
我们的六维 Token DSL 将控制信号提升至\textbf{人类可读的符号语言层面}。
这不仅是条件表示粒度的提升——频谱图需要从音频中提取，而我们的 Token 可从 MIDI 文件直接自动生成——
更关键的是，它为音乐生成的提示词工程建立了一套可扩展、可组合的标准化词汇体系。

% ─────────────────────────────────────────────────────────
\subsection{音乐结构化描述与段落级控制}

SegTune\cite{segtune}（ACL 2026 Oral）是该方向最具代表性的工作。
该方法允许用户为歌曲的不同段落（如 intro、verse、chorus）指定独立的自然语言提示词，
通过时间广播机制将段落级条件注入 MusicGen 的解码过程。
其核心贡献在于首次实现了音乐生成的\textbf{时序分段控制}——用户可以为 intro 描述
“轻柔的钢琴独奏”、为 chorus 描述“爆发式的管弦齐鸣”，全局 prompt 维持风格一致性。

然而，SegTune 的段落控制仍以\textbf{自然语言}为唯一描述形式。
“轻柔的钢琴独奏”是一种主观描述——模型的生成结果依赖于其对“轻柔”这一模糊概念的理解程度。
我们的六维 Token 体系则将“轻柔”精确化为 \texttt{[DYN:p] [ART:legato] [TEX:sparse] [VOICES:2]}，
将“爆发”精确化为 \texttt{[DYN:f] [ART:marcato] [TEX:dense] [VOICES:12]}，
消除了自然语言中的语义模糊性。

\textbf{与本文的差异：}SegTune 的段落条件与我们的 Token 条件处于不同抽象层级——
它们分别处理\textbf{“用什么语气说”}（自然语言）和\textbf{“说什么内容”}（结构化参数）。
两种条件在本质上互补：SegTune 可替代我们的 Text Prompt 模块，
而我们的 Token DSL 可显著扩展 SegTune 的控制精度。
此外，我们引入的 MoE 热插拔架构在风格层面的可扩展性是 SegTune 等单模型方案所不具备的。

% ═══════════════════════════════════════════════════════════
% 第 5 章：结论与未来工作
% ═══════════════════════════════════════════════════════════